% Nêu nội dung thực hiện và phương pháp dự kiến để đánh giá hệ thống

\subsection{Nội dung 1: Khảo sát và phân tích bài toán dự đoán lưu lượng mạng và Traffic Engineering}

\begin{itemize}
    \item \textbf{Mục tiêu:}

    Nghiên cứu tổng quan về bài toán dự đoán lưu lượng mạng và tối ưu Traffic Engineering trong các mạng hiện đại. Khảo sát các phương pháp học máy, Federated Learning và Reinforcement Learning đã được áp dụng trong lĩnh vực này nhằm xác định hướng tiếp cận phù hợp cho đề tài.

    \item \textbf{Phương pháp thực hiện:}

    Thu thập và phân tích các bài báo khoa học, tài liệu chuyên ngành và các nghiên cứu gần đây liên quan đến Traffic Engineering, dự đoán lưu lượng, Federated Learning và Reinforcement Learning. Khảo sát các bộ dữ liệu mạng công khai, phân tích đặc điểm dữ liệu và các chỉ số đánh giá, từ đó lựa chọn mô hình, thuật toán và kiến trúc hệ thống phù hợp cho quá trình nghiên cứu.
\end{itemize}

\subsection{Nội dung 2: Thiết kế và xây dựng hệ thống dự đoán lưu lượng kết hợp tối ưu Traffic Engineering}

\begin{itemize}
    \item \textbf{Mục tiêu:}

    Thiết kế hệ thống có khả năng dự đoán lưu lượng mạng bằng Federated Learning và tối ưu Traffic Engineering bằng Reinforcement Learning. Hệ thống đảm bảo các thành phần hoạt động đồng bộ và có thể mở rộng khi số lượng nút mạng tăng lên.

    \item \textbf{Phương pháp thực hiện:}

    Xây dựng kiến trúc hệ thống gồm các nút hoặc miền mạng đóng vai trò FL Client để thu thập dữ liệu và huấn luyện mô hình cục bộ, FL Server để tổng hợp các mô hình bằng FedAdam, mô-đun dự đoán lưu lượng và tác tử RL để thực hiện Traffic Engineering. Kết quả dự đoán lưu lượng được kết hợp với trạng thái mạng hiện tại và sử dụng làm đầu vào cho tác tử RL để đưa ra quyết định lựa chọn tuyến. Quyết định này được chuyển đến SDN Controller để thực thi trên mạng.
\end{itemize}

\subsection{Nội dung 3: Triển khai hệ thống mô phỏng và xây dựng chương trình demo}

\begin{itemize}
    \item \textbf{Mục tiêu:}

    Xây dựng hệ thống mô phỏng hoàn chỉnh nhằm minh họa quá trình dự đoán lưu lượng và tối ưu Traffic Engineering theo thời gian thực. Hệ thống cho phép quan sát trực quan sự thay đổi của mạng trước và sau khi áp dụng phương pháp đề xuất.

    \item \textbf{Phương pháp thực hiện:}

   Triển khai môi trường mô phỏng SDN bằng Mininet với topo mạng NSFNET và bộ điều khiển Ryu. Xây dựng các luồng lưu lượng, thu thập trạng thái mạng và tích hợp mô hình Federated Learning cùng Reinforcement Learning để thực hiện điều khiển lưu lượng chủ động. Phát triển chương trình demo gồm giao diện hiển thị trạng thái mạng, biểu đồ lưu lượng theo thời gian, mức sử dụng các liên kết và tuyến truyền đang được sử dụng.

    Quy trình hoạt động của hệ thống demo gồm các bước:
    \begin{enumerate}
        \item Khởi tạo topo mạng và các luồng truyền dữ liệu.
        \item Thu thập thông tin lưu lượng trên từng liên kết.
        \item Dự đoán lưu lượng trong khoảng thời gian tiếp theo bằng mô hình Federated Learning.
        \item Đưa kết quả dự đoán vào tác tử Reinforcement Learning để lựa chọn tuyến truyền tối ưu.
        \item Cập nhật đường đi của lưu lượng trên mạng.
        \item Hiển thị kết quả trực quan thông qua dashboard gồm topology mạng, lưu lượng trên từng liên kết, mức sử dụng băng thông, tuyến truyền trước và sau tối ưu, thông lượng, độ trễ và tỷ lệ mất gói.
    \end{enumerate}
\end{itemize}

\subsection{Nội dung 4: Thực nghiệm và đánh giá hệ thống}

\begin{itemize}
    \item \textbf{Mục tiêu:}

    Đánh giá hiệu quả của hệ thống đề xuất thông qua nhiều kịch bản lưu lượng mạng khác nhau và so sánh với các phương pháp hiện có.

    \item \textbf{Phương pháp thực hiện:}

    Xây dựng các kịch bản thực nghiệm gồm lưu lượng bình thường, lưu lượng tăng đột biến, hiện tượng tắc nghẽn trên một hoặc nhiều liên kết và sự thay đổi lưu lượng theo thời gian. Thực hiện so sánh giữa phương pháp đề xuất với các phương pháp định tuyến truyền thống hoặc các mô hình không sử dụng Federated Learning.

    Các chỉ số đánh giá bao gồm:
    \begin{itemize}
    \item Sai số dự đoán lưu lượng: MAE(Mean Absolute Error), RMSE(Root Mean Square Error), MAPE(Mean Absolute Percentage Error).
    
    \item Hiệu quả Traffic Engineering: thông lượng (Throughput), độ trễ đầu cuối (End-to-End Delay), tỷ lệ mất gói (Packet Loss Ratio), tỷ lệ sử dụng liên kết (Link Utilization) và mức sử dụng liên kết lớn nhất (Maximum Link Utilization - MLU).

    \item Khả năng cân bằng tải giữa các liên kết và khả năng giảm tắc nghẽn của mạng.

    \item Thời gian hội tụ của mô hình và thời gian đưa ra quyết định định tuyến.

    \item Khả năng thích ứng của hệ thống trước sự thay đổi của lưu lượng mạng.
\end{itemize}
\end{itemize}

\subsubsection{Giới hạn của đề tài}

Đề tài tập trung xây dựng và đánh giá hệ thống trong môi trường mô phỏng, chưa triển khai trên hạ tầng mạng của nhà cung cấp dịch vụ thực tế. Quy mô topo mạng, số lượng nút và lưu lượng được giới hạn nhằm phù hợp với điều kiện tính toán và thời gian thực hiện khóa luận. Nghiên cứu tập trung đánh giá hiệu quả của việc kết hợp Federated Learning và Reinforcement Learning trong bài toán dự đoán lưu lượng và tối ưu Traffic Engineering, chưa xem xét các vấn đề như triển khai trên mạng SDN quy mô lớn, bảo mật mô hình hoặc tối ưu đa miền.